\documentclass[11pt]{article}
\usepackage{amsmath,amssymb,amsfonts}
\usepackage{geometry}
\geometry{margin=1in}
\usepackage{hyperref}
\usepackage{bm}

\title{Multiplicity-Embedded Vlasov Dynamics: \\
Prime-Indexed Recursive Corrections in an Energy-Conserving DG Framework}
\author{Citizen Gardens \\ The Foundation of Multiplicity}
\date{April 26, 2026}

\begin{document}
\maketitle

\begin{abstract}
This document presents a defensive publication for a prime-indexed, recursively
defined multiplicity framework embedded into classical kinetic plasma theory.
Stage~1 focuses on a 1D--1V Vlasov--Poisson / Vlasov--Amp\`ere system, extended
by a hierarchy of prime-indexed modes and a structured recursive correction
operator. The correction is implemented as a divergence-form modification of
the spatial flux in an energy-conserving discontinuous Galerkin (DG) scheme.
The report provides: (i) an executive summary of the conceptual and practical
novelty; (ii) a comprehensive mathematical specification of the multiplicity
lift, reconstruction, and recursive correction; (iii) a fully coupled
algorithmic description compatible with energy-conserving DG solvers; and (iv)
representative code skeletons illustrating how to integrate the framework into
existing DG Vlasov solvers.
\end{abstract}

\tableofcontents
\newpage
\section{Executive Summary}

\subsection{Conceptual core}

Multiplicity Theory, as deployed here, treats classical fields and
distributions not as monolithic objects but as \emph{reconstructed aggregates}
of prime-indexed, recursively coupled modes. The Stage~1 embodiment of this
idea uses a kinetic plasma model---the 1D--1V Vlasov system coupled to Poisson
or Amp\`ere equations---as the base, and augments it by:
\begin{itemize}
  \item introducing a family of \emph{prime-indexed} distribution modes
        $f_p^{(k)}(x,v,t)$, where $p$ is a prime and $k$ is a recursion depth;
  \item defining a reconstruction map
        \(
          f_0(x,v,t) = \sum_{p,k} w_{p,k} f_p^{(k)}(x,v,t)
        \),
        which recovers the observable distribution $f_0$;
  \item designing a recursive correction operator
        $\mathcal{Q}_p^{(k)}$ that couples each mode only to modes of
        strictly lower recursion order and/or prime index;
  \item embedding one such correction, 
        $\mathcal{Q}_2^{(1)} = \eta_{2,1}\,\partial_x(\rho_2^{(0)} f_2^{(1)})$,
        as an additional divergence-form contribution to the DG spatial flux.
\end{itemize}

The \emph{field}---electric potential or field in the Vlasov--Poisson /
Vlasov--Amp\`ere system---couples only to the reconstructed distribution
$f_0$, not to individual modes. Thus the multiplicity structure lives entirely
inside the kinetic sector, while the field ``sees'' only the observable,
reconstructed, classical distribution.

\subsection{Practical novelty}

The key differentiators relative to standard Vlasov theory and DG
implementations are:
\begin{enumerate}
  \item A prime-indexed mode hierarchy $(p,k)$ with a \emph{recursion order}
        restriction: each $\mathcal{Q}_p^{(k)}$ depends only on modes
        $(q,l)$ with $(q,l) < (p,k)$ in lexicographic order on $(l,p)$.
  \item A reconstruction $f_0 = \sum_{p,k} w_{p,k} f_p^{(k)}$ 
        used as the \emph{exclusive} source for the field equations
        (Poisson or Amp\`ere), so that multiplicity is a hidden structure.
  \item A specific divergence-form correction 
        $\mathcal{Q}_2^{(1)} = \eta_{2,1}\,\partial_x(\rho_2^{(0)} f_2^{(1)})$
        that is designed to be implementable as a flux modification in an
        energy-conserving DG framework.
  \item An experimental protocol which:
        (i) recovers classical Landau damping and standard instabilities in the
        limit of zero multiplicity coupling,
        (ii) activates a correction-only regime (recursion with no direct
        contribution of additional modes to $f_0$),
        and (iii) finally allows selected modes to contribute to $f_0$,
        probing the physical impact of multiplicity on observables.
\end{enumerate}

The goal of this defensive publication is to fix this architecture---including
its prime-indexed recursion, reconstruction, and DG flux-based implementation
strategy---as prior art for future work in multiplicity-augmented kinetic and
field theories.

\section{Base Classical System: 1D--1V Vlasov--Poisson / Vlasov--Amp\`ere}

\subsection{Domain and classical distribution}

Let $x \in [0,L_x]$ be a periodic spatial coordinate and
$v \in [-V_{\max},V_{\max}]$ a truncated velocity coordinate.
The classical electron distribution function is
\[
  f_0(x,v,t) \ge 0,
\]
with total mass
\[
  M = \int_0^{L_x}\int_{-V_{\max}}^{V_{\max}} f_0(x,v,t)\,dv\,dx.
\]

\subsection{Vlasov equation}

The collisionless 1D--1V Vlasov equation reads
\begin{equation}
  \partial_t f_0 + v\,\partial_x f_0 + \frac{q}{m} E(x,t)\,\partial_v f_0 = 0,
  \label{eq:classical_vlasov}
\end{equation}
where $q$ is the charge, $m$ is the mass, and $E(x,t)$ is the self-consistent
electric field.

\subsection{Field equations}

Two closures are considered.

\paragraph{Vlasov--Poisson.}
Define the charge density
\[
  \rho_0(x,t) = \int_{-V_{\max}}^{V_{\max}} f_0(x,v,t)\,dv.
\]
The electrostatic potential $\phi(x,t)$ satisfies
\begin{equation}
  -\partial_x^2 \phi(x,t) = \rho_0(x,t) - \rho_{\mathrm{bg}},
  \label{eq:poisson}
\end{equation}
with $\rho_{\mathrm{bg}}$ a neutralizing background and periodic boundary
conditions, and
\[
  E(x,t) = -\partial_x \phi(x,t).
\]

\paragraph{Vlasov--Amp\`ere.}
Define the current density
\[
  J_0(x,t) = \int_{-V_{\max}}^{V_{\max}} v\,f_0(x,v,t)\,dv.
\]
The electric field evolves via
\begin{equation}
  \partial_t E(x,t) = -J_0(x,t),
  \label{eq:ampere}
\end{equation}
or a suitably scaled variant consistent with the DG energy-conserving scheme.

\section{Multiplicity Lift and Reconstruction}

\subsection{Prime-indexed modes and recursion}

Fix a finite set of primes
\[
  \mathbb{P} = \{p \in \mathbb{N} : p \text{ prime}, p \le p_{\max}\}
\]
and a maximum recursion depth $k_{\max}\in\mathbb{N}$.
For each $(p,k)$ with $p\in\mathbb{P}$ and $0\le k\le k_{\max}$ define a mode
distribution
\[
  f_p^{(k)}(x,v,t),
\]
which is not assumed to be nonnegative but is assumed to have bounded moments
and lie in the DG approximation space.

Define the lexicographic order on pairs $(p,k)$ by:
\[
  (q,l) < (p,k)
  \quad\text{if}\quad
  \big( l < k \big) \text{ or }
  \big( l = k \text{ and } q < p \big).
\]

\subsection{Reconstruction map}

Introduce weights $w_{p,k}$ and define the reconstructed distribution
\begin{equation}
  f_0(x,v,t) = \sum_{p\in\mathbb{P}}\sum_{k=0}^{k_{\max}} w_{p,k}\,
               f_p^{(k)}(x,v,t).
  \label{eq:reconstruction}
\end{equation}
Weights are scaled by a small parameter $\epsilon\ll 1$:
\[
  w_{p,k} = \lambda_{p,k}\,\epsilon,
  \quad
  \lambda_{2,0} = 1,
\]
so that the classical limit is recovered when $\epsilon\to 0$ and only mode
$(p,k)=(2,0)$ is active.

On a discrete grid, with $f_{p,i,j}^{(k)}$ the DG representation in cell $i$,
velocity cell $j$,
\[
  f_{0,i,j} = \sum_{p,k} w_{p,k} f_{p,i,j}^{(k)}.
\]

\subsection{Reconstructed observables}

For each mode $(p,k)$ define the density and current
\[
  \rho_p^{(k)}(x,t) = \int f_p^{(k)}(x,v,t)\,dv,
  \quad
  J_p^{(k)}(x,t) = \int v\,f_p^{(k)}(x,v,t)\,dv.
\]
The reconstructed density and current are then
\begin{equation}
  \rho_0(x,t) = \sum_{p,k} w_{p,k} \,\rho_p^{(k)}(x,t),
  \qquad
  J_0(x,t) = \sum_{p,k} w_{p,k} \,J_p^{(k)}(x,t).
  \label{eq:rho0_J0}
\end{equation}
These reconstructed quantities are used exclusively in the field equations
\eqref{eq:poisson}--\eqref{eq:ampere}.

\section{Recursive Correction Operator}

\subsection{General form}

Each mode satisfies a Vlasov-type equation with a correction operator
$\mathcal{Q}_p^{(k)}$:
\begin{equation}
  \partial_t f_p^{(k)} + v\,\partial_x f_p^{(k)} 
  + \frac{q}{m} E(x,t)\,\partial_v f_p^{(k)}
  = \mathcal{Q}_p^{(k)}[f],
  \label{eq:mode_vlasov_general}
\end{equation}
where $E(x,t)$ is computed from $f_0$ via \eqref{eq:poisson} or
\eqref{eq:ampere}.

The correction $ \mathcal{Q}_p^{(k)}[f] $ is required to satisfy:
\begin{enumerate}
  \item \textbf{Recursion:} it depends only on modes $(q,l)$ with $(q,l) < (p,k)$.
  \item \textbf{Divergence form:} it must be expressible as
        \[
          \mathcal{Q}_p^{(k)} = \partial_x A_{p,k}(x,v,t)
                                + \partial_v B_{p,k}(x,v,t),
        \]
        so that, with appropriate boundary conditions, it preserves mass and
        energy at the continuous level.
  \item \textbf{Small-coupling scaling:} it is proportional to the same small
        parameter $\epsilon$ that appears in the weights, via coefficients
        $\eta_{p,k} = \mu_{p,k}\,\epsilon$.
\end{enumerate}

\subsection{Stage~1 prototype: a single x-divergence correction}

For the first prototype, only a single nontrivial correction is activated:
\begin{equation}
  \mathcal{Q}_2^{(1)}[f]
  =
  \eta_{2,1}\,\partial_x\Big(\rho_2^{(0)}(x,t)\,f_2^{(1)}(x,v,t)\Big),
  \label{eq:Q21_def}
\end{equation}
where
\[
  \rho_2^{(0)}(x,t) = \int f_2^{(0)}(x,v,t)\,dv
\]
is the density of the lower mode $(2,0)$. All other $\mathcal{Q}_p^{(k)}$ may
be set to zero in the minimal implementation.

Equation \eqref{eq:Q21_def} is a pure $x$-divergence and fits naturally into
the $x$-advection step of a split DG method.

\section{DG Discretization and Flux Formulation}

\subsection{Spatial cells and test functions}

Partition $[0,L_x]$ into $N_x$ cells
$ I_i = [x_{i-1/2}, x_{i+1/2}]$, $i=1,\dots,N_x$, with cell centers
$x_i$. Let $\phi(x)$ be a test function from the DG space supported on
$ I_i$.

\subsection{X-advection substep for a mode}

Consider the $x$-advection + correction part for a mode $(p,k)$:
\begin{equation}
  \partial_t f_p^{(k)} + v\,\partial_x f_p^{(k)}
  =
  \mathcal{Q}_p^{(k)}[f].
  \label{eq:x_step_continuous}
\end{equation}
For the prototype, \eqref{eq:Q21_def} is used when $(p,k)=(2,1)$ and
$\mathcal{Q}_p^{(k)}=0$ otherwise.

Multiplying \eqref{eq:x_step_continuous} by $\phi$ and integrating over
$I_i$ yields
\[
  \int_{I_i} \partial_t f_p^{(k)} \phi\,dx
  +\int_{I_i} v\,\partial_x f_p^{(k)}\,\phi\,dx
  =
  \int_{I_i} \mathcal{Q}_p^{(k)}[f]\,\phi\,dx.
\]

Integrating by parts:
\begin{align*}
  \int_{I_i} \partial_t f_p^{(k)} \phi\,dx
  &+ \Big[ \widehat{(v f_p^{(k)})} \,\phi\Big]_{x_{i-1/2}}^{x_{i+1/2}}
  - \int_{I_i} v\,f_p^{(k)}\,\partial_x\phi\,dx \\
  &= \Big[ \widehat{(\eta_{p,k}\rho_{\text{lower}} f_p^{(k)})}\,\phi\Big]_{x_{i-1/2}}^{x_{i+1/2}}
  - \int_{I_i}\eta_{p,k}\,\rho_{\text{lower}}(x,t)\,f_p^{(k)}\,\partial_x\phi\,dx,
\end{align*}
where $\rho_{\text{lower}}$ is the appropriate lower-mode density, and 
$\widehat{(\cdot)}$ denotes the numerical flux at cell interfaces.

Define the total flux
\[
  \widehat{F}^{x}_{p,k}
  =
  \widehat{(v f_p^{(k)})}
  +
  \widehat{(\eta_{p,k}\rho_{\text{lower}} f_p^{(k)})}.
\]
The weak form becomes
\begin{equation}
  \int_{I_i} \partial_t f_p^{(k)} \phi\,dx
  +
  \Big[ \widehat{F}^{x}_{p,k}\,\phi\Big]_{x_{i-1/2}}^{x_{i+1/2}}
  -
  \int_{I_i} (v+\eta_{p,k}\rho_{\text{lower}}) f_p^{(k)}\,\partial_x\phi\,dx
  = 0.
  \label{eq:DG_weak_form}
\end{equation}

\subsection{Upwind numerical fluxes}

For the Vlasov term, use the standard upwind flux with speed $v$:
\[
  \widehat{(v f)} = v^+ f^- + v^- f^+,
  \quad
  v^\pm = \frac12\big(v\pm |v|\big),
\]
where $f^-$ and $f^+$ are left and right traces at the interface.

For the correction term, with speed 
\[
  a(x,t) = \eta_{p,k}\,\rho_{\text{lower}}(x,t),
\]
use the same upwind logic:
\[
  \widehat{(\eta_{p,k}\rho_{\text{lower}} f)}
  =
  a^+ f^- + a^- f^+,
  \quad
  a^\pm = \frac12\big(a\pm |a|\big).
\]
On a periodic domain, these fluxes preserve discrete conservation and are
compatible with the energy-conserving structure of the base DG scheme.

\subsection{V-advection substep}

The $v$-advection term,
\[
  \partial_t f_p^{(k)} + \frac{q}{m}E(x,t)\,\partial_v f_p^{(k)}=0,
\]
is discretized as in the underlying DG solver, with appropriate upwinding in
$v$ and boundary treatment (e.g.\ zero-flux) at $v = \pm V_{\max}$.

\section{Time Integration: Strang Splitting plus SSPRK3}

\subsection{Operators}

Define DG operators:
\begin{itemize}
  \item $\mathcal{A}_x[f_p^{(k)}, E; \Delta t]$:
        one DG step for the $x$-equation
        $\partial_t f_p^{(k)} + v\,\partial_x f_p^{(k)} = \mathcal{Q}_p^{(k)}[f]$
        with time step $\Delta t$.
  \item $\mathcal{A}_v[f_p^{(k)}, E; \Delta t]$:
        one DG step for the $v$-equation
        $\partial_t f_p^{(k)} + \frac{q}{m}E\,\partial_v f_p^{(k)}=0$
        with step $\Delta t$.
\end{itemize}

\subsection{Strang splitting for one time step}

For a time step $\Delta t$, a second-order Strang splitting for each mode is:
\begin{enumerate}
  \item X half-step:
    \[
      f_p^{(k),*} = \mathcal{A}_x\big[f_p^{(k),n}, E^n; \tfrac{\Delta t}{2}\big].
    \]
  \item Reconstruction and field update:
    \begin{itemize}
      \item Reconstruct $f_0^*$ from $\{f_p^{(k),*}\}$ via \eqref{eq:reconstruction}.
      \item Compute $\rho_0^*$ and $J_0^*$ via \eqref{eq:rho0_J0}.
      \item Solve Poisson \eqref{eq:poisson} or Amp\`ere \eqref{eq:ampere} to
            obtain $E^*$.
    \end{itemize}
  \item V full-step:
    \[
      f_p^{(k),**} = \mathcal{A}_v\big[f_p^{(k),*}, E^*; \Delta t\big].
    \]
  \item Reconstruction and field update again:
    \begin{itemize}
      \item Reconstruct $f_0^{**}$, compute $\rho_0^{**}$, $J_0^{**}$.
      \item Solve for $E^{**}$.
    \end{itemize}
  \item X half-step:
    \[
      f_p^{(k),n+1} = \mathcal{A}_x\big[f_p^{(k),**}, E^{**}; \tfrac{\Delta t}{2}\big].
    \]
\end{enumerate}

This Strang map can be used as the ``RHS flow'' inside a 3-stage SSPRK3
scheme by composing three such maps with appropriate convex combinations.

\section{Diagnostics and Experimental Protocol}

\subsection{Conserved quantities}

At time $t_n$, after reconstructing $f_0$ and $E$, compute:
\begin{itemize}
  \item Total mass:
    \[
      M(t_n) = \int_0^{L_x} \rho_0(x,t_n)\,dx.
    \]
  \item Kinetic energy:
    \[
      K(t_n) = \frac12 \int_0^{L_x}\int_{-V_{\max}}^{V_{\max}}
               v^2\,f_0(x,v,t_n)\,dv\,dx.
    \]
  \item Field energy:
    \[
      U_E(t_n) = \frac12 \int_0^{L_x} E(x,t_n)^2\,dx.
    \]
  \item Total energy:
    \[
      \mathcal{E}(t_n) = K(t_n) + U_E(t_n).
    \]
\end{itemize}
Monitor relative drifts
\[
  \frac{|M(t_n)-M(0)|}{M(0)},\quad
  \frac{|\mathcal{E}(t_n)-\mathcal{E}(0)|}{\mathcal{E}(0)},
\]
and also the minimum of $f_0$ to check positivity.

\subsection{Landau damping benchmark}

For a Landau damping test:
\begin{itemize}
  \item Initial condition: perturbed Maxwellian
    \[
      f_0(x,v,0) = \frac{n_0}{\sqrt{2\pi v_{th}^2}}
                   \exp\Big(-\frac{v^2}{2 v_{th}^2}\Big)
                   \big(1 + A\cos(k_0 x)\big).
    \]
  \item Classical baseline: only mode $(2,0)$ active, $\epsilon=0$,
        recover known Landau damping rate.
  \item Correction-only phase: modes $(2,0)$ and $(2,1)$ evolve with
        nonzero $\eta_{2,1}$, but $w_{2,1}=0$ so $f_0$ is unaffected
        directly by $(2,1)$; test conservation and any indirect effect
        through recursion.
  \item Full multiplicity phase: $w_{2,1}\neq 0$, small, so $(2,1)$
        contributes to $f_0$; study changes to damping rate and
        phase-space structure as functions of $\epsilon$.
\end{itemize}

\section{Representative Code Skeletons}

\subsection{Mode storage and reconstruction (Python-like pseudocode)}

\begin{verbatim}
# Mode storage
f_modes = {}       # key: (p, k), value: array of shape (Nx, Nv)
weights = {}       # w_{p,k}
eta = {}           # eta_{p,k} for corrections

# Initialize modes
f_modes[(2,0)] = init_maxwellian_perturbed()  # classical mode
f_modes[(2,1)] = np.zeros((Nx, Nv))           # correction mode

weights[(2,0)] = 1.0
weights[(2,1)] = 0.0    # initially does not contribute to f0

eta[(2,1)] = mu_21 * epsilon

def reconstruct_f0(f_modes, weights):
    f0 = np.zeros_like(next(iter(f_modes.values())))
    for (p,k), fpk in f_modes.items():
        f0 += weights[(p,k)] * fpk
    return f0
\end{verbatim}

\subsection{Lower-mode density and correction flux}

\begin{verbatim}
def compute_rho_mode(f_mode, v_grid):
    # integrate over v using quadrature
    rho = np.trapz(f_mode, v_grid, axis=1)  # shape (Nx,)
    return rho

def x_flux_with_correction(f_mode, v, rho_lower, eta_pk):
    # f_mode: shape (Nx, Nv)
    # v: velocity grid, shape (Nv,)
    # rho_lower: shape (Nx,)
    # compute upwind fluxes at cell interfaces

    # base Vlasov flux (v f)
    flux_vlasov = compute_upwind_flux_x(f_mode, v)

    # correction flux (eta * rho_lower * f)
    # need rho_lower at interfaces
    rho_edge = interpolate_to_edges(rho_lower)     # shape (Nx+1,)
    a_edge = eta_pk * rho_edge                     # advection speed at edges
    flux_corr = compute_upwind_flux_x(f_mode, a_edge)

    return flux_vlasov + flux_corr
\end{verbatim}

\subsection{Strang-split time step over all modes}

\begin{verbatim}
def time_step(f_modes, E, dt):
    # X half-step for all modes
    for (p,k) in f_modes:
        if (p,k) == (2,1):
            rho20 = compute_rho_mode(f_modes[(2,0)], v_grid)
            flux = x_flux_with_correction(f_modes[(p,k)], v_grid,
                                          rho20, eta[(2,1)])
        else:
            flux = compute_upwind_flux_x(f_modes[(p,k)], v_grid)
        f_modes[(p,k)] = dg_update_x(f_modes[(p,k)], flux, dt/2)

    # reconstruct f0 and update field
    f0 = reconstruct_f0(f_modes, weights)
    rho0 = compute_rho_mode(f0, v_grid)
    E = solve_field(rho0)  # Poisson or Ampère

    # V full-step for all modes
    for (p,k) in f_modes:
        f_modes[(p,k)] = dg_update_v(f_modes[(p,k)], E, dt)

    # reconstruct f0 and update field
    f0 = reconstruct_f0(f_modes, weights)
    rho0 = compute_rho_mode(f0, v_grid)
    E = solve_field(rho0)

    # X half-step
    for (p,k) in f_modes:
        if (p,k) == (2,1):
            rho20 = compute_rho_mode(f_modes[(2,0)], v_grid)
            flux = x_flux_with_correction(f_modes[(p,k)], v_grid,
                                          rho20, eta[(2,1)])
        else:
            flux = compute_upwind_flux_x(f_modes[(p,k)], v_grid)
        f_modes[(p,k)] = dg_update_x(f_modes[(p,k)], flux, dt/2)

    return f_modes, E
\end{verbatim}

\section{Conclusion}

This defensive publication fixes as prior art a specific architecture for
embedding a prime-indexed, recursively coupled multiplicity structure into a
classical 1D--1V Vlasov–Poisson / Vlasov–Amp\`ere system implemented with an
energy-conserving DG scheme. The essential elements are:
\begin{itemize}
  \item prime-indexed modes with recursion order $(p,k)$;
  \item reconstruction of the observable distribution $f_0$ as a weighted sum
        of modes;
  \item recursive correction operators $\mathcal{Q}_p^{(k)}$ that depend only
        on lower modes and are implemented as divergence-form flux
        modifications; and
  \item strict use of the reconstructed $f_0$ as the source for classical
        field equations, leaving the multiplicity structure hidden in the
        kinetic sector.
\end{itemize}
The combination of these features with an energy-conserving DG discretization
and a structured experimental protocol (baseline, correction-only,
full-multiplicity) constitutes a coherent and implementable framework that is
distinct from existing kinetic, DG, and quantum-corrected Vlasov literature.

\appendix
\section*{Mathematical Appendix}

\section{Functional Setting and Notation}

\subsection{Phase-space domain and norms}

Let $x \in [0,L_x]$ be the spatial variable with periodic boundary conditions,
and $v \in [-V_{\max}, V_{\max}]$ the velocity variable. The phase space is
\[
  \Omega := [0,L_x] \times [-V_{\max}, V_{\max}].
\]

We work with the standard Lebesgue spaces
\[
  L^2(\Omega) = \Big\{ g : \Omega \to \mathbb{R} \;\Big|\;
     \|g\|_{L^2(\Omega)}^2 := \int_0^{L_x}\int_{-V_{\max}}^{V_{\max}}
     |g(x,v)|^2\,dv\,dx < \infty \Big\},
\]
and similarly $L^2_x := L^2([0,L_x])$, $L^2_v := L^2([-V_{\max},V_{\max}])$.

For functions $g(x,v)$ we also use the mixed norms
\[
  \|g\|_{L^2_x L^2_v}
  :=
  \Big(
    \int_0^{L_x}\!\!\int_{-V_{\max}}^{V_{\max}} |g(x,v)|^2\,dv\,dx
  \Big)^{1/2},
  \qquad
  \|g\|_{L^\infty_x L^2_v}
  :=
  \operatorname*{ess\,sup}_{x\in[0,L_x]}
  \Big(
    \int_{-V_{\max}}^{V_{\max}} |g(x,v)|^2\,dv
  \Big)^{1/2}.
\]

Spatial derivatives are understood in the sense of periodic distributions in
$x$, and we consider Sobolev spaces $H^1_x$ with the usual norm
\[
  \|h\|_{H^1_x}^2
  :=
  \int_0^{L_x} \big( |h(x)|^2 + |\partial_x h(x)|^2 \big)\,dx.
\]

\subsection{Multiplicity modes}

For each prime $p$ and recursion depth $k$, we have a mode
\[
  f_p^{(k)} : \Omega \times [0,T] \to \mathbb{R},
\]
which we assume lies in $L^\infty(0,T; L^2(\Omega))$ and has bounded
first-order derivatives in $x$ in the sense that
\[
  \partial_x f_p^{(k)} \in L^2(0,T; L^2(\Omega)).
\]
These conditions are sufficient to justify the integration by parts steps used
below.

The lower-mode charges are defined by
\[
  \rho_p^{(k)}(x,t) := \int_{-V_{\max}}^{V_{\max}} f_p^{(k)}(x,v,t)\,dv,
\]
and we tacitly assume $\rho_p^{(k)} \in L^\infty(0,T;H^1_x)$ when needed.

\section{Continuous Energy Balance with Multiplicity Correction}

In this section we analyze the continuous (non-discrete) Vlasov dynamics for a
single mode with correction
\[
  \mathcal{Q}_2^{(1)} = \eta_{2,1}\,\partial_x\big(\rho_2^{(0)} f_2^{(1)}\big),
\]
and show that under suitable regularity and boundary conditions, this
correction preserves the total kinetic $+$ field energy of the reconstructed
distribution.

\subsection{Classical kinetic and field energy}

Let $f_0(x,v,t)$ be the reconstructed distribution
\[
  f_0(x,v,t) = \sum_{p,k} w_{p,k}\,f_p^{(k)}(x,v,t).
\]
The kinetic energy is
\[
  K(t) = \frac12 \int_{\Omega} v^2 f_0(x,v,t)\,dv\,dx,
\]
and for the electrostatic case we define the field energy as
\[
  U_E(t) = \frac12 \int_0^{L_x} E(x,t)^2\,dx,
\]
with $E = -\partial_x \phi$ where $\phi$ solves
\[
  -\partial_x^2 \phi = \rho_0 - \rho_{\mathrm{bg}},
  \qquad
  \rho_0 = \int f_0\,dv.
\]

We recall the classical identity: for the Vlasov--Poisson system with
periodic boundary conditions, the total energy
\[
  \mathcal{E}(t) := K(t) + U_E(t)
\]
is exactly conserved in time.

\subsection{Conservation for a single mode with correction}

Consider a single mode $f := f_2^{(1)}$ satisfying
\begin{equation}
  \partial_t f + v\,\partial_x f
  + \frac{q}{m} E\,\partial_v f
  = \eta\,\partial_x\big(\rho(x,t) f \big),
  \label{eq:f_single_correction}
\end{equation}
where $\rho(x,t) := \rho_2^{(0)}(x,t)$ is a given lower-mode density and
$\eta := \eta_{2,1}$ is a real constant. For clarity we suppress the $(2,1)$
indices.

\subsubsection{Mass conservation}

Let
\[
  M_f(t) := \int_{\Omega} f(x,v,t)\,dv\,dx.
\]
Integrating \eqref{eq:f_single_correction} over $\Omega$, and using periodic
boundary conditions in $x$ and appropriate decay / zero-flux conditions in
$v$, we obtain:
\begin{align*}
  \frac{d}{dt} M_f(t)
  &= \int_{\Omega} \partial_t f\,dv\,dx \\
  &= -\int_{\Omega} \Big(
       v\,\partial_x f + \frac{q}{m} E\,\partial_v f
     \Big)\,dv\,dx
     + \eta \int_{\Omega} \partial_x(\rho f)\,dv\,dx.
\end{align*}
The integration over $v$ of $\partial_v f$ vanishes under the assumed
boundary conditions at $v = \pm V_{\max}$. The spatial integral of
$\partial_x(\rho f)$ vanishes due to periodicity. Hence
\[
  \frac{d}{dt} M_f(t) = 0.
\]

\subsubsection{Kinetic energy contribution}

Define the kinetic energy of $f$,
\[
  K_f(t) := \frac12 \int_{\Omega} v^2 f(x,v,t)\,dv\,dx.
\]
Multiplying \eqref{eq:f_single_correction} by $\frac12 v^2$ and integrating
over $\Omega$, we have
\begin{align*}
  \frac{d}{dt} K_f(t)
  &= \frac12 \int_{\Omega} v^2 \partial_t f\,dv\,dx \\
  &= -\frac12 \int_{\Omega} v^3\,\partial_x f\,dv\,dx
     -\frac12 \frac{q}{m} \int_{\Omega} v^2 E\,\partial_v f\,dv\,dx
     +\frac12 \eta \int_{\Omega} v^2\,\partial_x(\rho f)\,dv\,dx.
\end{align*}
The first and third terms can be integrated by parts in $x$. With periodic
boundary conditions:
\[
  \int_{\Omega} v^3\,\partial_x f\,dv\,dx
  = \int_{-V_{\max}}^{V_{\max}} v^3 
    \left( \int_0^{L_x} \partial_x f\,dx \right)\,dv = 0,
\]
\[
  \int_{\Omega} v^2\,\partial_x(\rho f)\,dv\,dx
  = \int_{-V_{\max}}^{V_{\max}} v^2
    \left( \int_0^{L_x} \partial_x(\rho f)\,dx \right)\,dv = 0.
\]
Thus,
\[
  \frac{d}{dt} K_f(t)
  = -\frac12 \frac{q}{m} \int_{\Omega} v^2 E\,\partial_v f\,dv\,dx.
\]
Integrating by parts in $v$ and using the fact that $v^2 E$ is independent of
$v$,
\begin{align*}
  \int_{\Omega} v^2 E\,\partial_v f\,dv\,dx
  &= \int_0^{L_x} E(x,t)
     \left(\int_{-V_{\max}}^{V_{\max}} v^2 \partial_v f\,dv \right)\,dx \\
  &= \int_0^{L_x} E(x,t)
     \Big( v^2 f(x,v,t) \big|_{-V_{\max}}^{V_{\max}}
           - \int_{-V_{\max}}^{V_{\max}} 2v\,f(x,v,t)\,dv
     \Big)\,dx.
\end{align*}
If $f$ decays sufficiently fast or if we impose zero flux at $v=\pm V_{\max}$,
the boundary term vanishes, and we obtain
\[
  \int_{\Omega} v^2 E\,\partial_v f\,dv\,dx
  = -2 \int_0^{L_x} E(x,t)
                  \left(\int_{-V_{\max}}^{V_{\max}} v\,f(x,v,t)\,dv\right)dx.
\]
Denote the mode current
\[
  J_f(x,t) := \int_{-V_{\max}}^{V_{\max}} v\,f(x,v,t)\,dv.
\]
Then
\[
  \frac{d}{dt} K_f(t)
  = \frac{q}{m} \int_0^{L_x} E(x,t)\,J_f(x,t)\,dx.
\]

\subsubsection{Field energy contribution}

If $E$ evolves according to the full reconstructed current $J_0$ via
Amp\`ere-type dynamics,
\[
  \partial_t E(x,t) = - J_0(x,t),
\]
and $J_0$ includes the contribution from $J_f$ (weighted by $w_{2,1}$ in the
reconstruction), then a standard calculation shows that the total contribution
from all modes,
\[
  \mathcal{E}(t) = K_0(t) + U_E(t),
\]
where $K_0$ is the kinetic energy of $f_0$, satisfies
\[
  \frac{d}{dt} \mathcal{E}(t) = 0.
\]
Crucially, the correction term $\eta\,\partial_x(\rho f)$ does not introduce
additional bulk production or dissipation in $K_f$, because its contribution
vanishes after integration by parts in $x$ due to periodicity. Thus, the
continuous energy balance of the reconstructed system is preserved.

\subsection{Extension to the full mode hierarchy}

The above reasoning extends to the full family of modes $\{f_p^{(k)}\}$, under
two conditions:
\begin{itemize}
  \item each correction $\mathcal{Q}_p^{(k)}$ is of divergence form in $x$ and
        $v$ with suitable boundary conditions; and
  \item the field $E$ is sourced only by the reconstructed observables
        $(\rho_0, J_0)$, not by individual modes separately.
\end{itemize}
In that case, the total energy associated with $f_0$ and $E$ satisfies the
same conservation law as in the classical Vlasov theory.

\section{Operator Norm Bounds for the Recursive Correction}

In this section we treat the multiplicity correction
\[
  \mathcal{Q}[f] := \eta_{2,1}\,\partial_x(\rho_2^{(0)} f_2^{(1)})
\]
as an operator acting on the mode $f_2^{(1)}$. We derive a basic operator norm
bound in $L^2(\Omega)$.

\subsection{Assumptions}

Let $\rho(x,t) := \rho_2^{(0)}(x,t)$ be given. Assume that for a fixed time
$t$:
\begin{itemize}
  \item $\rho \in W^{1,\infty}([0,L_x])$, so that
        $\|\rho\|_{L^\infty_x}$ and $\|\partial_x \rho\|_{L^\infty_x}$ are
        finite;
  \item $f := f_2^{(1)}(\cdot,\cdot,t) \in H^1_x L^2_v$, i.e.\ $f$ and
        $\partial_x f$ are in $L^2(\Omega)$.
\end{itemize}
We define the operator $\mathcal{Q}$ at time $t$ by
\[
  (\mathcal{Q} f)(x,v) := \eta\,\partial_x(\rho(x) f(x,v)).
\]

\subsection{Pointwise bound}

Using the product rule,
\begin{equation}
  (\mathcal{Q} f)(x,v)
  = \eta \big( (\partial_x \rho(x)) f(x,v)
              + \rho(x)\,\partial_x f(x,v)\big).
  \label{eq:Q_operator}
\end{equation}
Hence, pointwise in $(x,v)$,
\[
  |(\mathcal{Q} f)(x,v)|
  \le |\eta|\Big( \|\partial_x \rho\|_{L^\infty_x}\,|f(x,v)|
                  + \|\rho\|_{L^\infty_x}
                    \,|\partial_x f(x,v)|\Big).
\]

\subsection{$L^2$ bound}

Squaring and integrating over $\Omega$,
\begin{align*}
  \|\mathcal{Q} f\|_{L^2(\Omega)}
  &\le |\eta| \Big[
    \|\partial_x \rho\|_{L^\infty_x}\,\|f\|_{L^2(\Omega)}
    + \|\rho\|_{L^\infty_x}\,\|\partial_x f\|_{L^2(\Omega)}
  \Big].
\end{align*}
Defining the mode-dependent norm
\[
  \|f\|_{H^1_x L^2_v}
  := \big( \|f\|_{L^2(\Omega)}^2
          + \|\partial_x f\|_{L^2(\Omega)}^2 \big)^{1/2},
\]
we can write
\begin{equation}
  \|\mathcal{Q} f\|_{L^2(\Omega)}
  \le |\eta|\,C_\rho\,\|f\|_{H^1_x L^2_v},
  \label{eq:Q_bound_H1}
\end{equation}
where
\[
  C_\rho := \max\{\|\partial_x \rho\|_{L^\infty_x}, \|\rho\|_{L^\infty_x}\}.
\]

Thus $\mathcal{Q}$ is bounded as an operator
\[
  \mathcal{Q}: H^1_x L^2_v \to L^2(\Omega),
\]
with operator norm
\[
  \|\mathcal{Q}\|_{H^1_x L^2_v \to L^2(\Omega)}
  \le |\eta|\,C_\rho.
\]

\subsection{Perturbative regime}

The small-coupling regime is defined by $|\eta|\ll 1$. If we regard the
classical transport operator
\[
  \mathcal{L} f := - v\,\partial_x f - \frac{q}{m} E\,\partial_v f
\]
as the dominant part and $\mathcal{Q}$ as a perturbation, then
\eqref{eq:Q_bound_H1} ensures that for sufficiently small $|\eta|$ the
combined operator
\[
  \mathcal{L}_\eta := \mathcal{L} + \mathcal{Q}
\]
is a bounded perturbation of $\mathcal{L}$ when viewed as an operator on
$H^1_x L^2_v$. Standard semigroup theory then implies stability and
well-posedness results under the classical hypotheses on $\mathcal{L}$,
with perturbations controlled by $|\eta| C_\rho$.

\section{Discrete Operator Considerations}

Although the main rigorous bounds above are written at the continuous level,
it is important to emphasize how the structure of $\mathcal{Q}$ interacts with
a DG discretization.

\subsection{DG projection of the correction operator}

Let $V_h$ be the DG finite element space in $x$, and $P_h$ the $L^2$
projection onto $V_h$ (for each fixed $v$). The discrete correction operator
acts on a DG mode $f_h \in V_h \otimes L^2_v$ as
\[
  \mathcal{Q}_h f_h := P_h\big[\eta\,\partial_x(\rho_h f_h)\big],
\]
where $\rho_h$ is the DG projection of the lower-mode density.

Because $P_h$ has operator norm 1 in $L^2(\Omega)$ and $\partial_x$ commutes
with $P_h$ on $V_h$, we obtain the analogous bound
\[
  \|\mathcal{Q}_h f_h\|_{L^2(\Omega)}
  \le |\eta|\,C_{\rho_h}\,\|f_h\|_{H^1_x L^2_v}
\]
with $C_{\rho_h}$ defined as in the continuous case but with $\rho$ replaced
by $\rho_h$.

\subsection{Discrete energy arguments}

The energy-conserving DG schemes for Vlasov--Poisson / Vlasov--Amp\`ere are
built on flux formulations which, for the classical transport operator
$\mathcal{L}$, lead to an exact discrete energy identity. Because the
multiplicity correction enters only as an additional DG-consistent flux term
with the same upwind structure and pure divergence in $x$, the discrete energy
arguments extend directly: the DG inner product of $\mathcal{Q}_h f_h$ with
$v^2$ produces only interface terms that cancel under periodic boundary
conditions, exactly as in the continuous case.

Thus, at the discrete level,
\[
  \mathcal{E}_h(t) = K_{0,h}(t) + U_{E,h}(t)
\]
is preserved modulo time discretization errors and the usual DG quadrature
approximations, with no extra artificial dissipation or growth induced by the
multiplicity correction.

\section{Summary of Analytical Guarantees}

The explicit bounds and identities established in this Appendix show that:
\begin{itemize}
  \item The multiplicity correction $\mathcal{Q}_2^{(1)}$ preserves the
        continuous mass and total energy of the reconstructed Vlasov system
        under standard boundary conditions.
  \item As an operator, $\mathcal{Q}_2^{(1)}$ is bounded from
        $H^1_x L^2_v$ to $L^2(\Omega)$, with norm controlled by
        $|\eta|$ and $\|\rho_2^{(0)}\|_{W^{1,\infty}_x}$.
  \item In a DG approximation, the projected correction operator
        $\mathcal{Q}_h$ obeys analogous bounds and can be incorporated as an
        additional flux term without breaking the discrete energy structure of
        an energy-conserving DG Vlasov solver.
\end{itemize}
These facts justify treating the Stage~1 multiplicity correction as a small,
energy-neutral perturbation of the classical Vlasov dynamics, both at the
PDE level and in the numerical scheme.

\section*{Addendum: Fokker--Planck Contributions to Energy Bounds}

In this addendum we state how the inclusion of a Fokker--Planck (FP) collision
operator modifies the continuous energy balance and the corresponding discrete
DG energy inequality, using the notation of the Mathematical Appendix.

\subsection*{A. Continuous Fokker--Planck energy balance}

Let $f_0(x,v,t)$ be the reconstructed distribution as in the Appendix, and let
the kinetic equation for $f_0$ be augmented by a Fokker--Planck operator
$C_{\mathrm{FP}}[f_0]$ in velocity:
\begin{equation}
  \partial_t f_0 + v\,\partial_x f_0
  + \frac{q}{m}E(x,t)\,\partial_v f_0
  = C_{\mathrm{FP}}[f_0],
  \label{eq:FP_vlasov_addendum}
\end{equation}
with $E$ obtained from $\rho_0$ as before. For definiteness, consider the
standard one-dimensional Fokker--Planck operator
\begin{equation}
  C_{\mathrm{FP}}[f_0]
  :=
  \nu\,\partial_v\big( v f_0 + \theta\,\partial_v f_0 \big),
  \label{eq:FP_operator_addendum}
\end{equation}
where $\nu>0$ is a collision frequency and $\theta>0$ the equilibrium
temperature parameter. Assume vanishing fluxes at $v=\pm V_{\max}$ so that
boundary integrals in $v$ vanish.

Define the kinetic and field energies
\[
  K(t) := \frac12\int_{\Omega} v^2 f_0(x,v,t)\,dv\,dx,
  \qquad
  U_E(t) := \frac12\int_0^{L_x} E(x,t)^2\,dx,
\]
and the total energy $\mathcal{E}(t) := K(t)+U_E(t)$ as in the Appendix. Then,
proceeding as in the collisionless case but with \eqref{eq:FP_vlasov_addendum}
and \eqref{eq:FP_operator_addendum}, one finds
\begin{equation}
  \frac{d}{dt}\mathcal{E}(t)
  = -D_{\mathrm{FP}}(t),
  \label{eq:FP_energy_balance_cont}
\end{equation}
where the Fokker--Planck dissipation functional $D_{\mathrm{FP}}(t)$ is
non-negative. For the model operator \eqref{eq:FP_operator_addendum}, a
standard calculation yields
\begin{equation}
  D_{\mathrm{FP}}(t)
  =
  \nu \int_{\Omega}
    \Big(
      \theta\,|\partial_v f_0(x,v,t)|^2
      + v^2\,|f_0(x,v,t)|^2
    \Big)\,dv\,dx
  \;\ge\; 0.
  \label{eq:FP_dissipation_cont}
\end{equation}
Thus, in the presence of Fokker--Planck collisions, the total energy is no
longer conserved but decays monotonically according to
\eqref{eq:FP_energy_balance_cont}.

Importantly, the multiplicity correction operators discussed in the main text
and Appendix remain pure divergence terms in $x$, so they do not contribute to
$D_{\mathrm{FP}}(t)$ and remain \emph{energy-neutral} at the continuous level.

\subsection*{B. Discrete DG Fokker--Planck energy inequality}

In the DG discretization, denote by $f_{0,h}(x,v,t)$ the DG approximation of
$f_0$, and by $C_{\mathrm{FP},h}[f_{0,h}]$ the DG discretization of the FP
operator, obtained by writing $C_{\mathrm{FP}}[f_0]$ in conservative form in
$v$ and applying DG fluxes at velocity interfaces.

Let
\[
  K_h(t) := \frac12\int_{\Omega} v^2 f_{0,h}(x,v,t)\,dv\,dx,
  \qquad
  U_{E,h}(t) := \frac12\int_0^{L_x} E_h(x,t)^2\,dx,
\]
and $\mathcal{E}_h(t) := K_h(t)+U_{E,h}(t)$ be the discrete kinetic, field,
and total energies, where $E_h$ is obtained from the reconstructed $\rho_{0,h}$
through the same DG Poisson or Amp\`ere solver as in the collisionless scheme.

Assuming:
\begin{itemize}
  \item the DG transport fluxes in $(x,v)$ are chosen as in the energy-conserving
        scheme of the Appendix, so that the collisionless discrete transport is
        skew-symmetric with respect to the discrete energy inner product;
  \item the DG FP operator $C_{\mathrm{FP},h}$ is constructed with symmetric
        (or interior-penalty) diffusive fluxes for the $\theta\,\partial_v^2 f$
        part and consistent upwind fluxes for the $\partial_v(v f)$ part,
\end{itemize}
one obtains a discrete analogue of \eqref{eq:FP_energy_balance_cont}:
\begin{equation}
  \frac{d}{dt}\mathcal{E}_h(t)
  = -D_{\mathrm{FP},h}(t),
  \label{eq:FP_energy_balance_discrete}
\end{equation}
with a discrete dissipation functional $D_{\mathrm{FP},h}(t)$ satisfying
\begin{equation}
  D_{\mathrm{FP},h}(t)
  \;\ge\; 0
  \quad\text{for all DG states } f_{0,h},
  \label{eq:FP_Dh_nonnegative}
\end{equation}
typically expressible as a sum of non-negative cell and interface contributions
(velocity gradients and jumps of $f_{0,h}$).

Thus, the collisionless multiplicity-embedded DG scheme of the main text and
Appendix, which satisfies a discrete conservation law
\[
  \frac{d}{dt}\mathcal{E}_h(t) = 0,
\]
is upgraded in the presence of a DG Fokker--Planck operator to the discrete
energy inequality \eqref{eq:FP_energy_balance_discrete}--\eqref{eq:FP_Dh_nonnegative}.
The multiplicity corrections, remaining pure $x$-divergence fluxes, do not alter
this inequality and remain energy-neutral at the discrete level.
\nocite{*}
\bibliographystyle{unsrt}  
\bibliography{references}  


\end{document}
